% DELTA23 characters workstream. Author: GPT-6 Astra Pro (/root/characters).
% First proof reading completed on 2026-11-21 Europe/Berlin.
% This file uses the input notation G=GT_0-hat, H=ker(lambda).

\subsection{A full relative character from the harmonic descent defect}

\paragraph{D23.C1 (proved from the cited inputs).}
Let $G=\widehat{\mathrm{GT}}_0$, $H=\ker\lambda$, and let
$D=D_P:G\to G$ be the endomorphism of
DELTA\_ON\_GT A.14, in its literal current marking.  Write
$a_2=\epsilon_2W$ and $a_3=\epsilon_3T$ on $H$.
The character
\[
 d_H(h)=a_2(h)+a_2(Dh)\quad(h\in H)
\]
has a unique continuous extension
\[
 d:G\longrightarrow\mathbb F_2,
 \qquad d|_{G_{\mathbb Q}}=0.
\]
For $g\in G$, its choice-independent formula is
\[
 d(g)=d_H(g\sigma^{-1}),\qquad
 \sigma\in G_{\mathbb Q},\quad\chi(\sigma)=\lambda(g).
\]
On $H$, this is the mixed duplication defect $d_2$ of
DELTA\_ON\_GT A.103 and A.107.  Its nonvanishing is not asserted.

\emph{Proof.}
The two individual characters $a_2,a_3$ are invariant under
conjugation by all of $G$: this is the separate normalizer
calculation in LM7\_PUSH E.65, equations (E.65.16)--(E.65.17),
not just invariance of their sum.  Since $D$ is a homomorphism
preserving $\lambda$, it maps $H$ to $H$, and
\[
 a_2(D(ghg^{-1}))=a_2(Dg\,Dh\,(Dg)^{-1})=a_2(Dh).
\]
Thus $d_H$ is a continuous, $G$-invariant character of $H$.
The arithmetic marking comparison DELTA\_ON\_GT A.62 says
$D(\tau)=c\tau c$ for $\tau\in G_{\mathbb Q}$, where
$c=(-1,1)$.  If also $\tau\in H$, conjugation invariance gives
$a_2(D\tau)=a_2(\tau)$; hence $d_H(\tau)=0$.

Every element of $G$ has the form $h\sigma$ with $h\in H$ and
$\sigma\in G_{\mathbb Q}$, because the cyclotomic character is
surjective.  If $h\sigma=h'\sigma'$, then
$h'^{-1}h=\sigma'\sigma^{-1}\in H\cap G_{\mathbb Q}$,
on which $d_H$ is zero.  Therefore $d(h\sigma)=d_H(h)$ is
well defined.  Multiplication in this form gives
\[
 (h\sigma)(k\tau)=h(\sigma k\sigma^{-1})(\sigma\tau),
\]
so conjugation invariance proves additivity.  The map
$H\times G_{\mathbb Q}\to G$, $(h,\sigma)\mapsto h\sigma$,
is a continuous surjection from a compact space to a Hausdorff
space.  The continuous function $(h,\sigma)\mapsto d_H(h)$ is
constant on its fibres, so the resulting $d$ is continuous.
The same decomposition proves uniqueness.  A.107, equation
(MD.8b), identifies its restriction with $d_2$.  No lifting
argument or pentagon relation enters this proof.\hfill$\square$

\paragraph{D23.C2 (proved; exact restricted classification).}
For an integer $n\geq0$, define the finite parity map
\[
 p_n:H\longrightarrow(\mathbb F_2^2)^{n+1},\qquad
 p_n(h)=\bigl((a_2(D^jh),a_3(D^jh))\bigr)_{0\leq j\leq n}.
\]
Among the continuous characters $G\to\mathbb F_2$ trivial on
$G_{\mathbb Q}$ whose restriction to $H$ factors through $p_n$,
all are precisely the span of
\[
 \Delta\circ D^j\quad(0\leq j\leq n),\qquad
 d\circ D^j\quad(0\leq j<n).
\]
This gives generators, not a claim that they are independent.
In particular, allowing finite strings of descendant frame parities
gives the $\mathbb F_2[D^*]$-module generated by $\Delta$ and $d$,
where $D^*\eta=\eta\circ D$.  It does not classify characters
without this factorization hypothesis.

\emph{Proof.}
Let $S_n=p_n(H)$, an $\mathbb F_2$-subspace of the displayed
finite vector space.  By DELTA\_ON\_GT B.14 the pair $(a_2,a_3)$
on $G_{\mathbb Q}\cap H$ has image $\{(0,0),(1,1)\}$.
The arithmetic formula $D\tau=c\tau c$, together with separate
conjugation invariance of the two parities, shows that the image
of $G_{\mathbb Q}\cap H$ under $p_n$ is exactly the diagonal line
spanned by the vector with all $2n+2$ entries equal to one.

A character factoring through $p_n$ is a linear functional on
$S_n$.  Extend it to the ambient vector space by extending a
basis of $S_n$ to a basis of that space.  Since the diagonal
vector lies in $S_n$ and its functional value is zero, the
extended functional has an even sum of coordinate coefficients.
The even-sum coefficient vectors are spanned by the following
$2n+1$ edge vectors: the sums of the two coordinates at a fixed
level $j$, and the sums of the harmonic coordinates at consecutive
levels $j,j+1$.  To check this without any rank assumption on
$S_n$, view the $2n+2$ coordinate positions as vertices: these
edges form a tree.  Removing a leaf and its incident edge eliminates
one specified coefficient; induction leaves the zero coefficient
at the final vertex because the total sum is even.  The corresponding
functions on $H$ are respectively
\[
 (a_2+a_3)\circ D^j=\Delta|_H\circ D^j,
 \qquad (a_2+a_2\circ D)\circ D^j=d_H\circ D^j.
\]
They extend to the global relative characters in the statement by
E.65 and D23.C1.  Two global characters trivial on $G_{\mathbb Q}$
that agree on $H$ agree on $G=HG_{\mathbb Q}$.  This proves both
containments and the asserted exact restricted classification.
\hfill$\square$

\paragraph{D23.C3 (proved; a separation and the harmonic lifting obstruction).}
Assume the actual element $k$ of DELTA\_ON\_GT LI.8
(integrated statement A.94) with
$Dk=1$, $\rho_k=0$, $q_{12}(\beta_k)=1$, and $\Delta(k)=1$.
Then
\[
 \Delta(k)=1,\qquad
 (\Delta\circ D^j)(k)=0\ (j\geq1),\qquad
 (d\circ D^j)(k)=0\ (j\geq0).
\]
In particular, if $d$ is nonzero it is independent of $\Delta$.
More generally $\Delta$ is outside the span of every positive
iterate of $\Delta$ and every iterate of $d$.

The elementary coprime-centralizer construction used in E.58
cannot, with identity relative action and a seed lying in the
relative kernel, be turned into a witness with
$q_{12}(\beta)=W=\pm S$ at three by replacing its prime two
component with a prime three component.

\emph{Proof.}
The retained frame and zero shear in LI.8 give $a_2(k)=0$.
Since $Dk=1$, the definition in D23.C1 yields $d(k)=0$ and
all positive iterates vanish.  Evaluating any proposed linear
relation at $k$ proves the separation claims.

For the last assertion, the relative lift in E.58 first chooses an
open normal subgroup $U$ and uses a seed whose finite frame value
belongs to $q(U)$, with $q(U)$ a $p$-group.  A harmonic value
$\pm S\in\mathrm{SL}_2(\mathbb F_3)$ has order four, so such
a $q(U)$ must be a two-group: it cannot be a three-group or any
odd-primary group.  But the torsion element being centralized is
$s$, of order two.  The hypothesis that its order be prime to the
pro-$p$ kernel, required for the coprime complement-conjugacy
argument, fails at this forced prime $p=2$.  Therefore that specific
construction does not establish the required centralizer lift.
This proves an obstruction to the proposed proof, not a theorem
that no other harmonic lift exists.\hfill$\square$
